\documentclass[10pt,twocolumn]{article}

\usepackage[T1]{fontenc}
\usepackage{mathpazo}
\usepackage[letterpaper,margin=0.85in,columnsep=0.3in]{geometry}
\usepackage{microtype}
\usepackage{amsmath}
\usepackage{amssymb}
\usepackage{booktabs}
\usepackage{array}
\usepackage{enumitem}
\usepackage[dvipsnames,table]{xcolor}
\usepackage[colorlinks=true,linkcolor=black,urlcolor=NavyBlue,citecolor=NavyBlue]{hyperref}

\definecolor{yaleblue}{HTML}{00356B}
\definecolor{diagtint}{HTML}{E9EEF5}
\definecolor{dnred}{HTML}{B3412C}

% interaction-figure helpers: own-base tint, mechanical marker, direction arrows
\newcommand{\own}[1]{\cellcolor{diagtint}\itshape #1}
\newcommand{\mech}{\textsuperscript{\,\textdegree}}
\newcommand{\up}{\textcolor{yaleblue}{\large$\blacktriangle$}\,}
\newcommand{\dn}{\textcolor{dnred}{\large$\blacktriangledown$}\,}
\newlength{\lblw}
\newlength{\colw}

\setcounter{secnumdepth}{1}
\usepackage[skip=0.4em plus 2pt]{parskip}

\newcommand{\tblgroup}[1]{\addlinespace[3pt]\multicolumn{4}{@{}l}{\textbf{\scriptsize\uppercase{#1}}}\\[1pt]}
\newcommand{\bibentry}{\par\hangindent=1.3em\noindent}

\title{\color{yaleblue}\textbf{An integrated model for scoring taxes at the top}}
\author{\normalsize The Budget Lab at Yale}
\date{}

\begin{document}
\maketitle

\section{The idea}

This memo describes the modeling enhancements we've made recently in service
of producing more-accurate and internally consistent scores of packages aimed
at taxing the rich. The basic idea is this:
\begin{enumerate}[leftmargin=1.4em,itemsep=0.25em]
\item Top-tax packages often combine ordinary-income, capital-gains, business,
estate, and/or wealth taxes in one proposal. \emph{Example: a proposal to
increase the corporate rate, increase capital gains rates, and tax unrealized
gains at death.}
\item These tax bases overlap both mechanically (e.g., a dollar of taxes paid
by a corporation reduces the shareholder tax base) and behaviorally (e.g.,
reducing a tax rate differential across two substitutable tax bases should
reduce avoidance opportunities). \emph{Example: the corporate tax rate hike
reduces after-tax profits which in turn reduces the tax base for the capital
gains tax, while at the same time the tax on gains at death reduces the
incentive to defer realizations, which raises the capital gains tax base.}
\item Some of these interactions have no empirical evidence to rely on, since
they arise from combinations of policies that have never been enacted.
\emph{Example: evidence on the realization elasticity is estimated solely
under a policy regime of step-up; we have no directly observable estimate of
how the elasticity changes if gains are taxed at death.}
\item Therefore, to the extent possible, we should build one integrated
microsimulation model with a ground-up, internally consistent representation
of the relationship between business incomes, household incomes, household
wealth, inheritance flows, etc., parameterized by behavioral rules grounded in
economic theory and calibrated to current-law empirical elasticity moments.
\end{enumerate}

This approach represents an improvement over what we've done in the past,
which is to build a set of standalone models and either ignore these
interactions entirely or crudely assume relationships that might not hold
under different policy regimes.

In what follows we describe how we model these mechanical and behavioral
interactions.

\section{How the model is put together}

This work is an extension of the Budget Lab's existing income and payroll tax
microsimulation. Some of it is new data and some is new modeling.

\subsection*{New data}

We train a nonparametric machine learning model on the SCF and the Forbes
billionaire list, where it estimates the joint distribution between
demographics, income, and wealth. We use this model to produce conditional
distributions of wealth and its composition for each record, then sample from
it, generating a wealth profile (including unrealized gains) for each tax
unit. We also fit annual mortality probabilities for each person, adjusted for
demographics and income. These new data are prerequisites for the integration
of tax bases: for each record, we can now calculate income and payroll taxes,
expected estate and deemed-realization taxes in the case of death, and wealth
taxes, and we can model adjustments to income and asset values when corporate
taxes change.

\subsection*{New modeling}

With this data in hand, we build several new modules in this simulator. These
modules are functions which map policy changes to tax base changes (behavioral
responses), or tax base changes to other tax base changes (mechanical
interactions). The modules apply in a fixed sequence on shared data; we do
not iterate them to a joint fixed point. Figure~\ref{fig:interactions} maps the interactions these
modules produce: each row is a tax increase, each column a base that moves in
response. Appendix A gives the derivation behind each module.

\begin{figure*}[p]
\caption{How the tax bases interact}
\label{fig:interactions}
\vspace{0.4em}
\footnotesize
\renewcommand{\arraystretch}{2.2}
\setlength{\tabcolsep}{4pt}
\setlength{\lblw}{1.9cm}
\setlength{\colw}{\dimexpr(\textwidth-\lblw-12\tabcolsep)/6\relax}
\begin{tabular}{@{}>{\raggedright\arraybackslash\bfseries}p{\lblw}*{6}{>{\raggedright\arraybackslash}p{\colw}}@{}}
\toprule
 & \multicolumn{6}{c}{\textbf{\ldots moves this base}} \\
\cmidrule(l){2-7}
\textbf{An increase in\ldots} & Ordinary income & Capital gains & Corporate & Estate & Gains at death & Wealth \\
\midrule
Ordinary rates &
\own{\dn Underreporting rises; giving responds} &
\up Compensation converts into deferred equity gains &
\up Income shifts into C-corporation form &
\dn Evaders' assets leave the reported estate &
\up Converted gains add to the stock later held at death &
\dn Evaders' assets leave the reported wealth base \\

Capital gains rates &
\up Conversion unwinds; pay returns as salary &
\own{\dn Realizations deferred; retiming around effective dates} &
\dn All-in corporate rate rises; income shifts to pass-through &
--- &
\up Deferred gains accumulate and arrive at death\mech &
--- \\

The corporate rate &
\dn Rates of return competed down; interest, rents, pass-through income fall\mech &
\dn Equity marked down; later realizations shrink\mech &
\own{\dn Income shifts toward pass-through form} &
\dn Marked-down portfolios shrink estates\mech &
\dn Unrealized gains absorb the markdown\mech &
\dn Marked-down portfolios shrink the base\mech \\

The estate tax &
--- &
\up Holding until death shelters less; realizations rise &
--- &
\own{\dn Reported estates shrink} &
\dn Gains realized earlier leave less held at death &
--- \\

The tax on gains at death &
\up Deferral worth less; conversion into gains unwinds &
\up Deferral no longer escapes tax; realizations rise during life &
\dn Retain-to-death shelter loses value; income shifts to pass-through &
\dn Gains tax due at death is deductible from the estate\mech &
\own{\dn Realizations pulled forward shrink the pool taxed at death} &
--- \\

The wealth tax &
\dn Concealed assets' income leaves the return &
\up Carrying a deferred gain costs wealth tax; realizations rise &
--- &
\dn Concealed assets leave the reported estate &
\dn Earlier realizations leave fewer gains at death &
\own{\dn Valuation discounts and concealment shrink reported wealth} \\
\midrule
\multicolumn{7}{@{}p{\dimexpr\textwidth-2\tabcolsep\relax}@{}}{%
\emph{And any of them:} \dn every during-life tax increase is partly financed
out of saving, and the forgone wealth compounds --- shrinking capital income,
wealth, and estate bases in later years.\mech} \\
\bottomrule
\end{tabular}

\vspace{0.6em}
{\scriptsize \up = the base grows; \dn = the base shrinks. Shaded = own-base
response. \mech\ = mechanical interaction; others are behavioral responses.
--- = no first-order channel in the model.}
\end{figure*}

\subsubsection*{\textit{Behavioral response:} Realizing or deferring gains}

The biggest modeling innovation in this work is our new treatment of capital
gains realization behavior: a structural model of the decision to realize or
defer, embedded in the microsimulation. Households are grouped into
representative cohorts by age, with each cohort's tax rates, mortality, and
estate-tax exposure weighted by holdings of unrealized gains, so the cohort
parameters reflect the wealthy households who actually hold the gains. Each
agent weighs the tax costs of selling now against the value of deferral, which
depends on the current capital gains rate, expected future rates, what happens
to the gain at death, and expected time until death. Wealth and estate taxes
also enter the value of deferral: an annual wealth tax falls on the deferred
liability every year it is carried, and because income tax paid at death
shrinks the taxable estate, the estate tax partly offsets the cost of taxing
gains at death. The decision is structured as a Bellman problem where each
dollar of unrealized gain carries a non-tax motive for selling (e.g.,
liquidity or rebalancing) and the holder sells when that motive outweighs the
tax wedge of realizing now rather than deferring. Behavior responds on two
margins---a permanent margin that sets the level of realizations, and a
short-run timing margin that lets a calibrated fraction of realizations shift
a year in either direction around an anticipated rate change---mirroring the
permanent/transitory distinction in the empirical literature. We calibrate the
model so that responsiveness matches these elasticities under current law. The
result, for a given policy change, is an implied yearly path of the change in realizations for each age group, distributed down to simulated tax
units.

This structure has three benefits over the old reduced-form approach where
we'd simply apply an elasticity to the baseline path of gains. The first is
endogenous behavior: both the level of realizations and the realization
elasticity itself change with the policy environment, with no additional
assumptions (such as judgmentally adjusting the elasticity for a change in
step-up). The second is accounting consistency: the model tracks the stock of
unrealized gains, so gains that go unrealized when taxes change do not
disappear---they resurface later as sales during life or gains held at death.
The third is a byproduct: an estimated tax benefit of deferral used in other
behavioral calculations, like the business entity choice (where the value of a
dollar in a C corp depends on the value of deferral).

\subsubsection*{\textit{Behavioral response:} Converting ordinary income into gains}

Owners of businesses and certain workers can structure labor compensation as
deferred equity gains rather than salary. The amount converted responds to
changes in the gap between the taxpayer's ordinary rate and the present value
of tax on a dollar entering the deferred-gain state. That present value comes
from the realization piece above, so any policy that makes deferral less
attractive makes this margin of conversion less attractive as well.

\subsubsection*{\textit{Behavioral response:} Moving between corporate and pass-through form}

Business income shifts between pass-through and C-corporation form with the
difference in their all-in tax rates. The C-corporation rate includes the
entity-level tax plus shareholder-level tax on distribution, with retained
earnings priced at the same deferral value as above. The value of sheltering
in corporate form therefore falls when gains are taxed at death.

\subsubsection*{\textit{Behavioral response:} Underreporting of income}

When marginal tax rates rise, some income leaves the tax return through
evasion. Following DeBacker, Heim, and Yuskavage (working paper presented at the National Tax Association meetings),
reported income responds to the net-of-tax rate where the IRS lacks
third-party information reporting, i.e., in self-employment, partnership and
S-corporation income, etc. The model adjusts these income forms when rates
change and removes a matching share of closely held assets from the reported
wealth and estate bases (assuming that income hidden from the income tax is
not then presented to the wealth tax or the estate tax).

\subsubsection*{\textit{Behavioral response:} Charitable giving}

Cash charitable giving responds to its tax price --- one minus the marginal
subsidy rate --- on the intensive margin. Giving of appreciated assets, whose
tax price also depends on the capital gains rate, is not yet modeled.

\subsubsection*{\textit{Behavioral response:} Wealth tax avoidance and evasion}

Evidence from other countries shows that reported wealth falls as the marginal
tax rate on wealth rises. We assume this response is smaller for marketable
assets than for closely held or other privately traded assets. We also assume
that this response comes from two different margins. One is valuation
discounts, which represent ways of lowering the assessed value of an asset
that remains visible to the IRS. The other is concealment, for example the use
of trust vehicles or offshore arrangements, which move the asset out of sight
entirely, taking its income out of this record's income tax base and its value
out of the potential estate. Marketable securities have market prices that
cannot be argued down, so their response is all concealment; the size of the
response from closely held assets is assumed to be split evenly between the
two components.

\subsubsection*{\textit{Behavioral response:} Estate tax avoidance and evasion}

The estate tax is computed record by record from the same balance sheets used
everywhere else, with valuation discounts, deductions, and spousal portability
modeled explicitly and reported estates benchmarked to SOI filings (so
avoidance already embedded in current-law estates is part of the baseline).
When a reform raises the estate tax, reported estates shrink further. This is
a pure reporting response (i.e., what heirs actually receive does not change),
and the tax is passed to heirs in the microdata by matching large estates to
large inheritances.

\subsubsection*{\textit{Mechanical interaction:} Tax-driven dissaving}

Taxes are paid out of either forgone consumption or forgone saving (assuming,
as we do in conventional scores, that real income is fixed). If some fraction
of new tax at time $t$ is paid out of saving --- that is, equivalently, if
MPCs are less than 100\% --- then at time $t{+}1$, all else equal, there
should be less private wealth (and lower estate value for those who die), less
capital income flowing from that wealth, and thus less in income tax. Until
this point, we implicitly assumed that all new taxes were paid out of
consumption, since capital income flows were fixed across all scenarios.

We now have the machinery to vary this parameter (the fraction of new tax
financed by lower saving) so that changes in taxes affect the path of wealth
going forward. We assume this parameter varies with wealth: near the bottom of
the distribution, higher taxes mostly displace consumption (i.e., high MPCs),
while at the top it is financed almost entirely out of saving (i.e., low
MPCs). The financed share then compounds: wealth that is not accumulated stops
earning returns, so the gap between baseline and policy wealth widens over
time. The consequences land in other bases years later: less capital income to
tax during life, a smaller wealth-tax base each year, and smaller estates at
death. A large during-life tax therefore partially cannibalizes wealth-based
revenue later in the budget window, and increasingly beyond it.

\subsubsection*{\textit{Mechanical interaction:} Corporate tax incidence}

A corporate tax rate increase reduces after-tax profits that get distributed
to shareholders. Equity markets should reprice down as stock values fall by
the PDV of lower profits. And corporate distributions (dividends and buybacks)
should also fall. Per Harberger logic, the burden does not all stay on
corporate equity: while profits from rents are simply capitalized into share
prices, the normal return to capital is competed down economy-wide as capital
reallocates, so the burden migrates to other assets like fixed income and
noncorporate business. For a given corporate tax rate increase, we quantify
the equity markdown, the reduction in corporate distributions, and the
reduction in the pretax returns for noncorporate capital. Then we adjust each
record's wealth and income based on portfolio composition. These changes flow
into the rest of the calculator, reducing income tax, capital gains, estate
taxes, wealth taxes, and so on.

\section{How each piece is parameterized}

Table~\ref{tab:params} lists the parameters, grouped by module. It also has
an ``epistemic status'' column with a rough, gut-level assessment of how
confident we are in each parameter value.

Most are taken directly from the
empirical literature but two are set differently. The first is $\eta$, the strength of the realization response in the realization model, calibrated so that the full model matches our preferred permanent realization elasticity of $-0.6$. The second is $\sigma$, the strength of ordinary-vs-gains income
conversion, set residually in a calibration exercise that targets a total elasticity of taxable income (ETI) of 0.25: we run the model with all other behavioral margins turned on and find
the value of $\sigma$ that delivers the target.

Note that this model produces \emph{conventional} estimates, so substantive
economic changes like labor supply or savings rate changes are excluded from
the behavioral adjustment margins. (The dissaving channel is a mechanical
description of how taxes are paid, not a model of how saving responds to rates
of return.)

\begin{table*}[p]
\caption{Parameters, by module}
\label{tab:params}
\small
\renewcommand{\arraystretch}{1.6}
\newcommand{\pform}[1]{\newline{\scriptsize\itshape #1}}
\begin{tabular}{@{}p{7.2cm}p{2.0cm}p{5.1cm}p{1.6cm}@{}}
\toprule
\textbf{Parameter} & \textbf{Value} & \textbf{Basis} & \textbf{Epistemic status} \\
\midrule
\tblgroup{Capital gains realization}
Long-run response ($\eta$)\pform{realizations scale by $\exp(-\eta\times{}$change in the tax cost of realizing now versus deferring$)$} & 2.48 & Calibrated on the full model to our preferred realization elasticity of $-0.6$ & High \\
\addlinespace[0.55em]
Share of realizations that can retime $\pm 1$ year\pform{the fraction of each year's realizations free to shift one year toward the lower-rate year} & 0.25 & Calibrated to short-run elasticity of $-1.2$ & High \\
\addlinespace[0.55em]
Valuation/compliance discount under deemed realization at death\pform{gains deemed realized at death are scaled down by this fraction before tax} & 25\% & Calibrated to JCT score & Low \\
\tblgroup{Income reporting and form}
Income conversion ($\sigma$)\pform{percent of the eligible compensation pool converted per point of change in the ordinary-versus-deferred-gain wedge} & 0.16 & Residual: set so the total top ETI is 0.25 (Saez, Slemrod, and Giertz 2012) & Low \\
\addlinespace[0.55em]
Entity-shifting semi-elasticity\pform{percent of pass-through business income changing form per point of change in the all-in rate differential} & 0.63 & Prisinzano and Pearce (2018) & High \\
\addlinespace[0.55em]
Underreporting elasticities: self-employment / partnership \& S-corp / rent\pform{percent change in reported income per one percent change in the net-of-tax rate} & 0.046 / 0.052 / 0.040 & DeBacker, Heim, and Yuskavage; zero under third-party reporting & Medium \\
\addlinespace[0.55em]
Charitable giving price elasticity (cash)\pform{percent change in cash giving per one percent change in the price of giving} & $-0.5$ & Literature central value & Low \\
\tblgroup{Wealth and estates}
Reported-wealth semi-elasticity: marketable / closely held\pform{reported wealth scales by $\exp($elasticity${}\times{}$marginal wealth-tax rate$)$, by asset class} & $-7$ / $-17$ & Budget Lab review of international evidence & Low \\
\addlinespace[0.55em]
Concealment share of the avoidance response: marketable / closely held\pform{the concealed share exits the income and estate bases too; the valuation share exits reported wealth only} & 100\% / 50\% & Marketable assets have observable prices; private ones do not & Medium \\
\addlinespace[0.55em]
Reported-estate net-of-tax elasticity\pform{the reported estate retains the net-of-tax ratio raised to this power} & 0.16 & Kopczuk and Slemrod (2001); range 0.10--0.22 & Medium \\
\tblgroup{Mechanical interactions}
Share of new tax paid out of saving, by wealth rank\pform{the share of each new tax dollar paid out of saving; the forgone wealth compounds thereafter} & 0.04 $\rightarrow$ 0.80 & Consumption responses by wealth (Straub 2019 and related) & Medium \\
\addlinespace[0.55em]
Corporate incidence constants (e.g., normal-return share 0.375)\pform{exposure and repricing constants in the incidence formulas of Appendix A.8} & various & Provisional placeholders pending direct measurement & Low \\
\bottomrule
\end{tabular}
\end{table*}

\clearpage
\section*{Appendix A. Derivations for Section 2}

The sections below follow the order of the modules in Section 2 (charitable
giving, a one-parameter response, is documented in Appendix B). Each is
self-contained: symbols are defined where they first appear and are local to
their section.

\subsection*{A.1 Realizing or deferring gains}

\paragraph{Cohorts and the stock of unrealized gains.}
The model operates on representative age cohorts (for joint returns, the older
spouse's age), followed forward to age 119 using SSA life tables beyond the
microdata's range. The state variable is the cohort's stock of unrealized
gains. Each year the stock is drawn down by realizations and by death, where
the death regime determines what happens to the decedent's gains: under
step-up they are extinguished, under carryover they pass to heir cohorts
(allocated by a dollar-weighted heir-age distribution from the SCF), and under
deemed realization they are taxed. The model tracks the change in each
cohort's stock relative to baseline, so a gain that goes unrealized when rates
rise does not disappear; it remains in the stock and resurfaces as a later
sale or as a gain held at death.

Because unrealized gains are heavily concentrated, every cohort-level
input---the marginal tax rate on gains, mortality, estate-tax exposure, the
wealth-tax carrying cost defined below---is weighted by each household's share
of the cohort's gains, so the cell parameters reflect the wealthy households
who actually hold the gains. Mortality in particular is the probability that
the marginal dollar's holder dies, well below the cohort's average. The
cohort-level response is then distributed back to tax units in proportion to
a blend of their realized gains and their unrealized holdings.

\paragraph{The agent's problem.}
Consider one dollar of unrealized gain. Its holder has some non-tax motive for
selling---liquidity, rebalancing, consumption---summarized by a reservation
benefit $b \geq 0$ drawn from an exponential distribution. The dollar is sold
when $b$ exceeds the tax cost of realizing now rather than deferring, so the
share of the stock realized in a year, $r$, is the survival function of that
distribution evaluated at the tax wedge. Equivalently, the cohort chooses $r$
subject to a realization cost $C(r)$ whose marginal is
\begin{equation*}
C'(r) = \frac{1}{\eta}\,\ln\!\left(\frac{r}{r_B}\right),
\end{equation*}
where $r_B$ is the cohort's baseline realization rate and $\eta$ governs
response strength. The cost is zero at the margin under baseline policy,
$C'(r_B) = 0$, so baseline behavior is reproduced exactly.

The tax cost of realizing now rather than deferring is
\begin{equation*}
MC \;=\; \tau \;+\; \beta\,(1-m)\,\big(W_{\text{next}} - h\big) \;+\; \beta\, m\, F,
\end{equation*}
where $\tau$ is the current marginal tax rate on long-term gains (measured
through the tax calculator, so it embeds surtaxes and interactions), $\beta$
is the real discount factor built from the 10-year Treasury rate deflated by
CPI inflation year by year, $m$ is mortality, $W_{\text{next}}$ is next year's
continuation value of an unrealized dollar (the same problem solved one year
and one age ahead), $h$ is the annual wealth-tax cost of carrying the deferred
position, and $F$ is the value of the tax treatment at death. Setting the
marginal non-tax benefit equal to the marginal cost gives a closed form for
the scenario realization rate:
\begin{equation*}
r \;=\; r_B \,\exp\!\big(-\eta\,(MC - MC_B)\big), \qquad \text{capped at } 1,
\end{equation*}
where $MC_B$ is the marginal cost under baseline law. The response is a
constant semi-elasticity in the full deferral wedge, not in the current-year
rate alone.

\paragraph{What enters the deferral wedge.}
Expected future rates enter through $W_{\text{next}}$, which is solved by
backward induction over ages and years, so a pre-announced rate change moves
behavior before it takes effect. The remaining terms give precise form to the statement in Section 2 that wealth and estate taxes also affect the value of
deferral.

The death value is
\begin{equation*}
F \;=\; (1 - c_\varphi)\,\tau\,(1 - e).
\end{equation*}
Here $c_\varphi$ is the share of the gain's tax burden the holder bears at
death: $0$ under step-up (death forgives the tax entirely), $1$ under deemed
realization (death triggers the full tax, so $F = 0$ and deferral buys nothing
at death), and an intermediate value under carryover, reflecting how much of
the heir's eventual liability the holder internalizes. The regime can differ
by asset class, and $c_\varphi$ is the gain-weighted mix; gains bequeathed to
charity escape tax under any regime and are netted out. The factor $(1 - e)$
is the estate-tax offset: capital-gains tax due at death is deductible from
the taxable estate, so each dollar of forgiven gains tax is worth only
$(1 - e)$, where $e$ is the household's marginal estate rate. Below the estate
exemption, $e = 0$ and the term is inert.

The carrying cost $h$ is the product of the household's marginal wealth-tax
rate and its capital-gains rate, averaged with gain weights. Realizing today
removes the tax from the wealth base, while deferring keeps the full pre-tax
dollar in the base, so each year of continued deferral costs the wealth tax on
the deferred liability. An annual wealth tax therefore raises realizations
even with no change in the capital-gains schedule.

The 25 percent valuation and compliance discount under deemed realization
applies to the revenue collected from gains taxed at death, not to the
holder's incentive: the deferral decision sees the full statutory burden.

\paragraph{The short-run timing margin.}
Retiming around an anticipated rate change is a separate overlay. A calibrated
fraction of each year's realizations may shift by up to one year toward the
adjacent year with the lowest rate, with the shifted share proportional to the
rate gap. Under a permanent, uniform rate change every year looks the same and
the overlay nets to zero, so the permanent margin and the timing margin can be
calibrated independently.

\paragraph{Calibration.}
We calibrate $\eta$ on the full model rather than on the cohort problem in
isolation, so that the target elasticity is hit after mortality, death
treatment, discounting, and the allocation to tax units are all accounted for.
We run the simulator under a permanent capital-gains rate increase, measure
the proportional realization response at a thirty-year horizon per unit change
in the rate, and repeat across a small grid of $\eta$ values. The measured
response is linear in $\eta$ through the origin, so the calibration is an
inversion of a fitted line. The target is our preferred permanent realization
elasticity of $-0.6$, expressed as a semi-elasticity by dividing by the
baseline top rate of about $0.238$; with a fitted slope of $1.0155$ this gives
$\eta = 2.4825$, which we quote as $2.48$. The timeable fraction, $0.25$, is
calibrated the same way to a short-run realization elasticity of $-1.2$ around
an announced rate change. The deemed-realization discount of 25 percent is set
so that the model's ten-year revenue from taxing gains at death matches the
JCT score of that policy.

\paragraph{The exported object: the value of deferral.}
The realization model also produces, for each cohort and year, the expected
present value of tax per dollar entering the deferred-gain state, which we
write $\tau_{eq}$. It is computed by a backward recursion over the scenario's
own realization path: in each future year a deferred dollar pays $\tau$ if
realized, pays the death-triggered tax net of the estate offset if its holder
dies under deemed realization, and pays the carrying cost $h$ if it survives
unrealized under a wealth tax, all discounted at $\beta$. Under current law
with step-up, $\tau_{eq}$ is well below the statutory rate; taxing gains at
death pushes it toward $\tau$. This object is the price of the deferred-gain
state used by the income-conversion and entity-form margins (Sections A.2 and
A.3): any policy that makes deferral less valuable raises $\tau_{eq}$ and
thereby discourages recharacterizing ordinary income as gains, with no
separate assumption required.

\subsection*{A.2 Converting ordinary income into gains}

Owner-managers can take part of their compensation as appreciation in a
business they control rather than as salary or current profit: below-market
salary at a closely held C corporation, sweat equity realized on eventual
sale, profits interests and carried interest, founder stock. Each of these
routes trades a dollar taxed this year at ordinary rates for a dollar of
unrealized gain taxed later, or never. We model them together as a single
conversion margin.

The price of the salary path is the taxpayer's own marginal rate on that
compensation source, $m_{i\ell}$, measured through the calculator (payroll tax
included) for each source $\ell$: own wages, spouse wages, and active
partnership, S-corporation, and sole proprietorship income. The price of the
equity path is $\tau_{eq}(a,t)$, the equivalent tax rate on a dollar entering
the deferred-gain state: the expected present value of tax that dollar
eventually generates through sales during life and its treatment at death,
computed by the realization model of Section A.1 for each age group $a$ and
year $t$. The wedge between the two paths is
$W_{i\ell} = m_{i\ell} - \tau_{eq}(a_i, t)$, and the behavioral forcing is the
reform-induced change in that wedge,
\begin{equation*}
\Delta W_{i\ell} \;=\; \Delta m_{i\ell} \;-\; \Delta\tau_{eq}(a_i, t),
\end{equation*}
so a higher top ordinary rate widens the wedge while a policy that taxes gains
at death raises $\tau_{eq}$ and narrows it. The amount converted is
\begin{equation*}
\Delta c_{i\ell} \;=\; \sigma \,\Delta W_{i\ell}\, P_{i\ell},
\end{equation*}
where the pool $P_{i\ell}$ is wages in full and 75 percent of active
pass-through income, the labor-content share estimated by Smith, Yagan, Zidar,
and Zwick (2019). The response is recomputed each year from that year's wedge,
runs in either direction, and is capped so no compensation source turns
negative. Conversion is limited to tax units with taxable income at or above
the top ordinary bracket threshold and some active business income; in 2026
this is roughly one million tax units holding about \$1.5 trillion of
compensation. Converted dollars leave wages and pass-through income
immediately, shrinking the income and payroll tax bases, and enter the age
cohort's stock of unrealized gains in the following year, where they are
realized and taxed according to the dynamics of Section A.1 and meet whatever
death regime the policy specifies; the later tax is attributed at the cohort
level rather than traced back to the converting record.

The response parameter $\sigma$ is set residually. The evidence that
disciplines this margin is the elasticity of taxable income at the top, for
which we adopt the Saez, Slemrod, and Giertz (2012) central value of 0.25. But
the other reporting margins already generate most of that response: with
evasion, entity shifting, and charitable giving turned on and conversion
turned off, a five-point increase in the top ordinary rate produces a measured
elasticity of top-bracket taxable income (excluding net capital gains) of
about 0.20. Adding an independently estimated conversion elasticity on top
would count the same behavior twice, since those estimates are themselves
drawn from total-response evidence. We therefore run the full model with every
other margin on and solve for the $\sigma$ that brings the total to 0.25.
Because the converted amount is linear in $\sigma$, two runs pin the answer:
$\sigma = 0.16$, meaning 0.16 percent of the pool converts per percentage
point of wedge change, confirmed by a full run at that value. The calibration
disciplines the product of $\sigma$ and the pool definition, so a narrower
pool would simply imply a proportionally larger $\sigma$ with the same
aggregate response. The value is conditional on the rest of the stack and is
re-derived whenever the other margins change.

\subsection*{A.3 Moving between corporate and pass-through form}

Business income moves between pass-through and C-corporation form in response
to the difference in their all-in tax rates, following Pearce and Prisinzano
(2018). On the corporate side, a dollar of profit bears the entity-level tax
and then a shareholder-level tax on what remains. We assume a share
$\alpha = 0.45$ of after-tax corporate earnings is paid out currently and
taxed at the shareholder's marginal rate on long-term gains, $\tau_g$
(dividends and gains are assumed to face the same rate). The retained share is
priced at $\tau_{eq}(a,t)$, the equivalent deferred-gain tax rate from Section
A.1: a retained dollar is an unrealized gain in the shareholder's hands, and
its tax cost is the expected present value of tax it generates through later
realization and its treatment at death. The all-in corporate rate is
\begin{equation*}
\tau^{C}_i \;=\; \tau_c \;+\; (1-\tau_c)\,\big[\,\alpha\,\tau_{g,i} + (1-\alpha)\,\tau_{eq}(a_i,t)\,\big],
\end{equation*}
where $\tau_c$ is the statutory corporate rate. The pass-through side is the
taxpayer's ordinary marginal rate on active pass-through income, $\tau^{P}_i$,
inclusive of self-employment tax and any pass-through deduction.

The share of business income that changes form responds to the change in the
differential under the reform:
\begin{equation*}
\Delta q_i \;=\; \varepsilon\,\Delta\big(\tau^{C}_i - \tau^{P}_i\big), \qquad \varepsilon = 0.63,
\end{equation*}
a semi-elasticity: 0.63 percent of the taxpayer's business income shifts per
percentage point change in the differential. The value comes from
Prisinzano and Pearce's preferred estimate, rescaled by the pass-through sector's 60
percent share of business income at which it was identified. The shifting base
is each record's positive net pass-through business income across
partnerships, S corporations, and sole proprietorships. When income shifts
toward corporate form, it leaves the pass-through base, pays the entity-level
tax, and returns to the shareholder as after-corporate-tax distributions, so
the score records both the corporate receipts and the shareholder-level tax on
what is paid out; dollars leaving one base arrive in the other. The
shareholder-level tax on retained earnings is booked as a current-year
adjustment to the record's capital gains priced at $\tau_{eq}$, rather than
by adding retained dollars to the deferred-gain stock.

Because $\tau_{eq}$ enters the all-in corporate rate, the value of corporate
form as a shelter depends on the capital-gains and death-tax regime, not just
on the entity rates. Under stepped-up basis, retained earnings held to death
escape shareholder tax entirely and $\tau_{eq}$ is low, so corporate form
shelters effectively. A policy that taxes unrealized gains at death raises
$\tau_{eq}$, and with it $\tau^{C}$, even when no corporate parameter changes;
the differential moves against corporate form and previously sheltered income
shifts back toward pass-through status, where it is taxed currently. The model
produces this unwinding automatically.

\subsection*{A.4 Underreporting of income}

When marginal rates rise, some income stops being reported. We follow
DeBacker, Heim, and Yuskavage, who estimate the elasticity of noncompliance
with respect to the net-of-tax rate from IRS random-audit data and treat it as
a component of the elasticity of taxable income. For each income type $j$ with
elasticity $e_j$, reported income under the reform is
\begin{equation*}
\tilde y_{ij} \;=\; y_{ij}\left[\,1 + e_j\!\left(\frac{1-m_{ij}}{1-m^0_{ij}} - 1\right)\right],
\end{equation*}
where $m_{ij}$ and $m^0_{ij}$ are the taxpayer's marginal rates on that income
under the reform and under the baseline. Only income the IRS cannot see
through third-party information returns responds: self-employment income
(with farm income treated the same way), at 0.046; partnership and
S-corporation income, at 0.052; and rent, at 0.040. Wages, interest, and
dividends are covered by information reporting and do not respond; audit data
show misreporting rates near one percent for wages against a large majority of
unmatched sole proprietorship income. Only positive income responds---the
overstatement of losses and deductions is not modeled---and because these
estimates come from detected noncompliance in random audits, which miss some
sophisticated evasion at the very top, they are best read as a floor.

Evaded income does not vanish; it accumulates in the evader's hands, largely
inside the closely held businesses whose income went unreported. The model
keeps the balance sheet consistent with the tax return: a taxpayer who hides
business income from the income tax does not then present the underlying
assets to the wealth tax or the estate tax. For each record we form the evaded
share of closely held income,
\begin{equation*}
u_i \;=\; \frac{\sum_{\ell}\big(y_{i\ell} - \tilde y_{i\ell}\big)}{\sum_{\ell} y_{i\ell}},
\end{equation*}
summing over the record's positive closely held income sources, and remove
that share of its closely held assets from the reported wealth-tax base and
from the reported estate at death. Where a wealth tax also induces concealment
(Section A.5), the two are combined as a union,
$1 - a_i = (1 - c_i)(1 - u_i)$, where $c_i$ is the concealed fraction from the
wealth-tax response and $a_i$ the combined share removed, so an asset hidden
through both routes is removed only once. The taxpayer's actual balance sheet
is untouched: the assets still earn income, still enter the saving and
realization dynamics, and still pass to heirs in full. What shrinks is only
what the tax authority observes, which is what each reported base is computed
from.

\subsection*{A.5 Wealth tax avoidance and evasion}

The wealth tax is assessed on net worth, assets less debts, under a graduated
schedule whose exemption operates as a zero-rate bottom bracket. The marginal
rate $\tau^w$ is therefore zero below the exemption, and the avoidance
response is confined to households above it.

Reported wealth falls with the marginal rate under a semi-elasticity
specification applied separately to two asset classes: marketable assets
(cash, listed equities, bonds, retirement accounts, life insurance, trusts,
and other financial assets) and closely held assets (private business
interests, homes, and other real estate and nonfinancial assets). The reported
value of a household's class-$j$ assets is
\begin{equation*}
\tilde{W}_j = W_j \, e^{\varepsilon_j \tau^w},
\end{equation*}
where $W_j$ is the true value and $\varepsilon_j$ is the class
semi-elasticity, $-7$ for marketable and $-17$ for closely held assets. The
avoided fraction $f_j = 1 - e^{\varepsilon_j \tau^w}$ is approximately
$-\varepsilon_j \tau^w$ at low rates: about 7 percent of marketable and 16
percent of closely held wealth at a 1 percent marginal rate, and about 19 and
40 percent at 3 percent. Debts are reported in full. These elasticities are
carried over from the Budget Lab's standalone wealth-tax model and were
reviewed and retained for this work; we treat them as central values and
recommend sensitivity bands around them in published estimates.

The response decomposes into a valuation component and a concealment
component, with the concealment share set by asset class: 100 percent for
marketable assets, whose exchange prices cannot be argued down, so avoiding
them means hiding them; 50 percent for closely held assets, where legal
valuation discounts do half the work. The two components exit different bases.
A valuation discount lowers assessed wealth only: the asset's income still
appears on the income tax return, and its full value still enters the reported
estate. A concealed dollar leaves all three reported bases at once: it exits
reported wealth, its income (interest, dividends, pass-through income, rent,
and the gains realized when it is sold) exits the income tax base, and its
value exits the reported estate at death. The one exception is distributions
from retirement accounts, which remain fully reported because those accounts
are subject to third-party information reporting. The same consistency runs in
the other direction: income already hidden through the underreporting response
of Section A.4 removes the matching share of closely held assets from the
reported wealth and estate bases, and the two channels are combined so that
overlap is counted once.

These are reporting responses, not real ones. The household's true balance
sheet is unchanged, and everything computed from it is untouched:
mortality-weighted estates, the financing of taxes out of wealth, and exposure
to the corporate channel all read the true values, and heirs inherit concealed
wealth in full. Finally, the elasticities implicitly absorb migration and
expatriation, so wealth tax revenue results are best read as an upper bound on
that account.

\subsection*{A.6 Estate tax avoidance and evasion}

The estate tax is computed record by record from the same modeled balance
sheet used everywhere else. Let $G$ denote a record's economic gross estate,
the sum of its asset values. Reported gross estate applies a valuation
mapping,
\begin{equation*}
\tilde{G} = G \cdot r \, \bigl[\, 1 + (\rho - 1)\, \theta \,\bigr],
\end{equation*}
where $\theta$ is the share of assets held in pass-through businesses, $r$ is
an overall valuation factor, and $\rho$ is an additional discount concentrated
in pass-through wealth, reflecting minority-interest and marketability
discounts. The calibrated values are $r = 0.95$ and $\rho = 0.61$, chosen so
that modeled taxable returns and tax match SOI filings across estate-size
classes in the latest observed death year (2022). Debts $B$ are subtracted
directly from the data; non-debt deductions, dominated by charitable and
residual marital bequests, enter as a size-varying fraction $\delta$ of
reported gross estimated from SOI filings, rising from about 7 percent below
\$10 million to about 36 percent above \$50 million. The distributable estate
is $D = \max(\tilde{G} - B - \delta \tilde{G},\, 0)$. The unified base adds
lifetime taxable gifts back at a rate $\gamma \approx 0.09$ of reported gross,
also estimated from filings, and subtracts the decedent's income tax due at
death when current law makes it deductible. Liability applies the graduated
tentative schedule $T(\cdot)$ with the unified credit taken at the exclusion
amount $E$: $L(E) = \max\bigl(T(\text{base}) - T(E),\, 0\bigr)$.

Spousal portability is handled as explicit branches. A married decedent with a
living spouse is modeled at the event that both die, with liability $L(2E)$.
For unmarried decedents, including widows and widowers, the calculator
computes liability twice, once at the basic exclusion and once with the
survivor's unused spousal exemption added, and weights the two by an estimated
probability that such an exemption is available; two complete calculations are
required because the credit introduces a kink that an expected exemption would
misprice near the threshold. All of this is conditional on death. Expected
revenue applies each household's mortality probability $m$ as a weight at
aggregation: $\mathbb{E}[\text{tax}] = w\, m \,[\, p\, L_{\text{port}} +
(1-p)\, L \,]$, where $w$ is the record weight and $p$ the portability
probability.

Reported estates also respond to the estate tax itself. Following Kopczuk and
Slemrod (2001), the elasticity of the reported estate with respect to the
net-of-tax rate is 0.16, with a plausible range of 0.10 to 0.22. We implement
the exact finite-change form: the retained share of the reported base is
$\bigl[(1-\tau_S)/(1-\tau_B)\bigr]^{0.16}$, where $\tau_B$ and $\tau_S$ are
the record's marginal estate rates under baseline and reform, both measured
through the calculator. The retained fraction is applied to reported gross
estate rather than the taxable base --- the same convention as the wealth
concealment --- which slightly overstates the response above the exemption,
within the published range. The response keys off the change in the rate, never
its level, because baseline-level avoidance is already embodied in the
calibrated valuation factors; it is symmetric under rate cuts, and it reduces
the taxable base without changing what heirs receive.

Unrealized gains at death follow the regime being scored, set per asset class:
under step-up the gain escapes income tax; under carryover it passes to heirs
and resurfaces in their later sales; under deemed realization it is taxed on
the decedent's final return, less a 25 percent valuation and compliance
discount calibrated to the JCT score, with the resulting tax deductible
against the estate. The regime feeds back into during-life realization
behavior through the deferral model of Section A.1: the value of holding a
gain until death is the income tax escaped, $\tau(1-e)$, where $e$ is the
marginal estate rate, since income tax paid at death would have shrunk the
taxable estate. Under current law this offset is about 18 percent of the
death-tax value for the top decile of gain holders, so estate-only reforms
move realizations during life.

Estate tax reaches heirs by rank matching. Decedent estates, weighted by death
probability and sorted by distributable value, and inheritances, sorted by
size, form two cumulative-dollar ladders that are matched from the top down;
each heir bears the average tax rate of the estates occupying the
corresponding dollar interval, and the allocated total equals expected estate
revenue each year. Allocating the tax in proportion to inheritance would be
wrong for a tax with a high threshold: nearly all estates owe nothing, and a
proportional rule would assign tax to small inheritances that come from exempt
estates. The identifying assumption is that larger inheritances come from
larger estates. Inheritance amounts are gross of estate tax, so they do not
vary with policy; only the tax borne does. Deemed-realization tax at death,
which has no threshold, is allocated in proportion to inheritance instead.

Timing: estate tax on calendar-year-$t$ decedents is due nine months after
death, so revenue is booked in fiscal year $t+1$.

\subsection*{A.7 Tax-driven dissaving}

A new tax bill has to be paid out of something. A household that owes more can
consume less, or it can save less; there is no third option. This section
describes the accounting we use to track the second case. Let $s$ denote the
share of an additional dollar of tax financed out of saving rather than
consumption. We assume $s$ rises with wealth: about 0.04 near
the bottom of the within-age wealth distribution and 0.80 at the top, with a modest age
pattern layered on top (lower for the young and for retirees of ordinary
means, highest for peak earners, and flat at the top of the wealth
distribution, where the elderly do not run down wealth with age).

The profile is anchored to the evidence on consumption responses to persistent
income changes. For a permanent change in after-tax income, consumption moves
by the elasticity of consumption with respect to permanent income times the
household's consumption-to-income ratio, so $s = 1 - \varepsilon \cdot (C/Y)$.
With $\varepsilon \approx 0.7$ (Straub 2019) and a top-1-percent
consumption-to-income ratio of roughly 0.3 (Mian, Straub, and Sufi 2021), the
top of the distribution finances about 80 cents of a marginal tax dollar out
of saving; near the bottom, where constrained households consume close to
their income, the share is about 10 cents. Saving rates that reach half of
income in the top percentile (Dynan, Skinner, and Zeldes 2004) are consistent
with this gradient.

The channel operates on cohorts: cells defined by age and by percentile of net
worth within age. (Records with zero or negative net worth are left
untouched.) For each cell we track $P_t$, the gap between baseline wealth and
wealth under the policy, in dollars. Each year we measure, record by record,
the extra tax the household owes under the policy relative to baseline during
life (income, payroll, and any wealth tax), net of any change in the
household's pre-tax income that the policy itself caused (relevant only for
corporate reforms; see Section A.8). Call the cell total of this net cash-flow
hit $F_t$. It is measured on a run with behavioral responses switched on but
this channel switched off, so it reflects how households respond to the law
but not the channel's own feedback. The gap then evolves as
\begin{gather*}
P_{t+1} \;=\; G_t\,P_t \;+\; s\,F_t,\\
G_t \;=\; (1 + r_t) \;-\; s\,(\tau\,y + \tau_w),
\end{gather*}
where $r_t$ is the nominal growth rate of income per capita, which we use as
the return at which cohort wealth compounds; $\tau$ is the marginal tax rate
on the bundle of capital income the missing wealth would have generated
(dividends, interest, realized gains, and a portion of pass-through income),
measured through the tax calculator and aggregated to the cell; $y$ is the
cell's capital income per dollar of gross wealth; and $\tau_w$ is the marginal
wealth-tax rate. The subtraction inside $G_t$ is the feedback term: wealth
that was never accumulated throws off no taxable income and owes no wealth
tax, and that forgone tax is itself a cash-flow relief financed at the same
share $s$.

The accumulated gap is then pushed back onto records. Each record receives a
share of its cell's $P_t$ in proportion to its net worth; its asset values,
unrealized gains and their basis, capital income flows, and retirement
distributions are scaled down accordingly, and net worth is recomputed.
Because the wealth-tax and estate calculators run on the adjusted balance
sheet, they collect less automatically: a smaller wealth-tax base each year
and, weighted by each record's mortality, a smaller estate at death. This is
how a large during-life tax partly cannibalizes wealth-based revenue later in
the window and increasingly beyond it.

Two statements of discipline. First, this is an accounting description of how
taxes are paid, not a model of how saving responds to rates of return; nothing
in it moves when after-tax returns change. Second, it enters only the
behavioral (conventional) estimate. The static score, taxes under the new law
with behavior held fixed, is never touched.

\subsection*{A.8 Corporate tax incidence}

The module is organized around a single shock. The input is the gross change
in corporate receipts, before any individual-tax
offset. Holding aggregate
pre-tax income fixed, a corporate tax increase reduces after-tax profits
dollar for dollar, so $w_t$ --- read directly from the scored receipts
change --- corresponds conceptually to $\Delta\tau_t\,\Pi_t$ for pre-tax
profit $\Pi_t$ and rate change $\Delta\tau_t$. Everything below traces this one reduction
through the model: as current income flows, as asset values, and finally as
its consequence for individual receipts.

Income flows first. We split $w_t$ into a supernormal
component, borne by corporate equity alone (returns to market power and ideas
have no supply margin and are simply capitalized), and a normal-return
component, which spreads across all capital as capital reallocates away from
the corporate sector. We set the normal-return share at $\nu = 0.375$, leaving
0.625 on supernormal returns, consistent with Treasury and Tax Policy Center
estimates that about 60 percent of the corporate base is supernormal. The
reallocation is not instant: the normal-return burden phases in with the
replacement of the capital stock, $\lambda_t = 1 - (1 - 0.057)^t$, roughly
half complete after twelve years. Corporate equity therefore bears
\begin{equation*}
h_t \;=\; (1-\nu)\,w_t \;+\; \nu\,w_t\,\big[(1-\lambda_t) + \lambda_t\,\kappa\big],
\end{equation*}
where $\kappa = 0.40$ is corporate equity's share of the economy-wide stock of
normal capital: even in the long run, the single after-tax return that clears
the capital market keeps that fraction of the normal burden on corporate
capital. The remainder, $\nu\,w_t\,\lambda_t\,(1-\kappa)$, falls on
noncorporate capital, so interest (phasing in as debt is rolled over at lower
rates), rents, and a slice of pass-through income decline too. On each record,
dividends scale with corporate equity's after-tax profit hit,
$\phi_t = -h_t/\pi_t$, where $\pi_t$ is aggregate after-tax profit; at enactment,
before reallocation has begun, this is the familiar
$-\Delta\tau/(1-\tau)$.

The same profit reduction is capitalized into asset values. The market value of corporate
equity falls by the present value of the future after-tax profit reductions
that remain priced into corporate shares, computed by backward recursion at a
constant nominal rate, the 10-year Treasury yield at enactment plus a
5-percentage-point equity premium, with a growing-perpetuity terminal value
beyond the projection horizon. Expressed as a proportional discount $\mu_t$,
the markdown is applied to directly held equity and to the equity share of
funds, trusts, and retirement balances. Unrealized gains absorb the full price
change, since basis does not move; later realizations shrink on both margins,
the volume of sales (which tracks after-tax profits through repurchases) and
the gain per dollar sold (which tracks the markdown). Retirement distributions
scale with their marked-down balances, and estates reprice automatically
because the estate calculator runs on the adjusted balance sheet. Under a
permanent rate change the markdown decays toward the supernormal component as
reallocation returns the normal burden to the whole capital stock; under a
temporary change it is largest at announcement and unwinds by expiry.

Last, the consequence for individual receipts. The declines in
taxable dividends, interest, rents, gains, and retirement distributions flow
through the individual tax calculator, and individual receipts fall by
whatever the microsimulation produces given who actually holds the affected
income. The offset is therefore an output of the model, not an assumed
fraction of the corporate estimate. Household exposure is limited: we
attribute about 85 percent of dividends and half of realized gains to
corporate equity, and roughly 40
percent of the total burden falls on foreign investors, nonprofits, and
defined-benefit plans that do not appear on household returns. We do not
force household records to absorb that share; a reconciliation report carries
it as an unallocated remainder. The after-tax income losses also enter the dissaving channel of Section A.7, so part of the burden drains into smaller estates
decades later.

As in Section A.7, all of this affects only the behavioral estimate; the
static score remains a clean statement of the law. The constants in this
section (the supernormal split, $\kappa$, the exposure shares, the equity
premium) are provisional placeholders pending direct measurement, and we
present corporate-rate results as such.

\section*{Appendix B. Sourcing for Table 1}

This section expands the one-line ``Basis'' column of Table 1. For each
parameter we give the underlying citation, describe how the published estimate
maps into the model's units, and, where the value is calibrated rather than
taken from the literature, describe the calibration and date it. Precise
pinned values are given here; Table 1 rounds them.

\medskip\noindent\textbf{\small CAPITAL GAINS REALIZATION}

\textbf{Long-run response ($\eta$).} The pinned value is 2.4825; Table 1
rounds to 2.48. $\eta$ is the semi-elasticity in the realization rule---each
cell's discretionary realization rate scales by
$\exp(-\eta \times \text{change in the marginal tax cost of realizing})$---and
because the whole gain pool responds, it is also the model's aggregate
long-run semi-elasticity. It is calibrated so that the full model reproduces
our preferred permanent realization elasticity of $-0.6$, which sits inside
the empirical range: Dowd, McClelland, and Muthitacharoen (2015) estimate a
persistent elasticity of $-0.72$ (with a within-year elasticity of $-1.2$) on
a large panel of returns, Agersnap and Zidar (2021) find ten-year revenue
elasticities of $-0.3$ to $-0.5$ with revenue-maximizing rates of 38 to 47
percent, and the JCT and CBO scoring conventions sit in the same neighborhood.
The calibration (July 2026): we ran a permanent increase in the top statutory
rate on gains to 25 percent through the full model at several trial values of
$\eta$, measured the realized long-run response at simulation year 30, and
found the measured response linear in $\eta$; inverting the fitted line at the
target semi-elasticity of $-0.6/0.238 = -2.52$ (the elasticity converted at
the 23.8 percent top federal rate on long-term gains) gives 2.4825. The fitted
line's coefficients are kept in the calibration records so that future
re-derivations are arithmetic.

\textbf{Share of realizations that can retime $\pm 1$ year.} Pinned at 0.2542;
Table 1 rounds to 0.25. This is the fraction of baseline realizations free to
shift across adjacent years toward the lower-rate year, and it carries the
model's short-run announcement response. It is calibrated (July 2026) to our
preferred short-run elasticity of $-1.2$, which is the within-year estimate in
Dowd, McClelland, and Muthitacharoen (2015). Because the retiming margin nets
to zero under a uniform permanent rate change, $\eta$ is pinned first and
independently; given $\eta$, the timing share is then found by a direct
root-finding search on the full model against the short-run moment (target
semi-elasticity $1.2/0.238 = 5.04$).

\textbf{Valuation and compliance discount under deemed realization at death.}
A measurement parameter, not a behavioral response: it scales down the gains
deemed realized at death for valuation games and noncompliance that revenue
estimators assume but the realization model itself does not produce. At the 25
percent discount, the model's ten-year revenue from a
deemed-realization-at-death regime comes in near \$695 billion, against
roughly \$600 billion in JCT scoring of comparable proposals; a discount of 25
to 33 percent is also consistent with the valuation discount our estate-tax
calibration implies for pass-through assets. We intend eventually to replace
this single number with the estate side's asset-class-specific reporting
factors.

\medskip\noindent\textbf{\small INCOME REPORTING AND FORM}

\textbf{Income conversion ($\sigma$).} $\sigma = 0.16$ governs how much
top-bracket labor-type compensation is restructured into unrealized equity
gains when the wedge between the ordinary rate and the equity path's expected
present-value tax rate widens: 0.16 percent of the eligible compensation pool
converts per percentage point of wedge. It is the model's one residual
parameter. With every other reporting margin active (realization, entity
shifting, underreporting, charity), we run a five-point top-rate increase
through the full model, measure the total elasticity of top taxable ordinary
income, and set $\sigma$ so that this lands on 0.25---the central reading of
Saez, Slemrod, and Giertz (2012), whose survey puts the best available
estimates between 0.12 and 0.40. The value was re-derived in July 2026 after
revisions to the entity-shifting and underreporting modules, yielding 0.157,
shipped as 0.16; a confirmation run measures the total top ETI at 0.2508.
Because $\sigma$ is residual, it is re-derived whenever any other margin in
the stack changes; a standing check in the calibration records enforces this.

\textbf{Entity-shifting semi-elasticity.} From a Penn Wharton Budget Model
working paper by Prisinzano and Pearce (2018), who estimate how the share of
business income accruing to corporations responds to the
corporate/pass-through rate differential. We take their preferred estimate of
0.3788 (their Table IV.B) and divide by the pass-through sector's roughly 60
percent share of business income, giving a semi-elasticity of pass-through
income of $0.3788/0.6 \approx 0.63$ per percentage point of rate differential.
One departure from the paper's framework: where the realization model is
active, the tax value of deferral on retained corporate earnings is priced by
that model's expected present-value tax rate for the owner's age and the
prevailing death regime, rather than by
the paper's fixed approximation. The paper's identification comes from small
historical movements in the rate differential, and it cautions against
extrapolating far beyond them --- a caveat that applies to the large reforms
scored here.

\textbf{Underreporting elasticities.} From DeBacker, Heim, and Yuskavage
(2025), presented at the National Tax Association annual meetings in November
2025, using IRS National Research Program random-audit data from 2006 to 2017.
They estimate the elasticity of noncompliance with respect to the net-of-tax
rate as a component of the overall ETI, and we apply their values as
net-of-tax-rate elasticities of reported income, which is the authors' own
usage in their revenue illustration. The central values come from their pooled
cross-section subsample estimates on federal rate variation: 0.046 for filers
with self-employment income (applied to Schedules C and F) and 0.052 for
filers with partnership income (applied to partnership and S-corporation
income). The rent value, 0.040, is their Kansas difference-in-differences
estimate for Schedule E filers, the weakest-identified of the three. Their
quasi-experimental designs give larger values---0.09 for Schedule C filers in
the Kansas reform and up to 0.26 for sole proprietors in the bunching
design---so the centrals sit at the conservative end of the paper's
0.02-to-0.26 range. Wages, interest, and dividends are assigned no response,
reflecting third-party information reporting; only positive income is scaled,
so the overstated-loss and deduction margins are not modeled (a documented
omission---they find itemizer elasticities of 0.07 to 0.23).

\textbf{Charitable giving price elasticity (cash).} A tax-price elasticity of
$-0.5$ applied to cash contributions on the intensive margin, where the price
of giving is one minus the marginal subsidy rate. The value matches Randolph
(1995), who separates permanent from transitory price responses in panel data
and finds a permanent price elasticity of about $-0.5$. Later panel estimates
run somewhat larger---Bakija and Heim (2011) find persistent elasticities of
roughly $-0.7$ to $-1$---so $-0.5$ is a conservative central value for a
persistent response. The appreciated-asset giving margin, whose price depends
on the capital gains rate, is a documented gap.

\medskip\noindent\textbf{\small WEALTH AND ESTATES}

\textbf{Reported-wealth semi-elasticities.} The values $-7$ for marketable and
$-17$ for closely held wealth are assumptions carried over from our standalone
wealth-tax model, reviewed and accepted in July 2026; they are not estimates
from a single study. Units: reported wealth in each class scales by
$\exp(\text{rate} \times e)$, so at a 3 percent marginal rate they imply
roughly 19 percent erosion of reported marketable wealth and 40 percent of
closely held. The quasi-experimental literature spans a wide range that
brackets these values: Seim (2017) finds small, mainly reporting responses in
Sweden; Jakobsen, Jakobsen, Kleven, and Zucman (2020) find sizable long-run
responses at the top in Denmark; Br\"ulhart, Gruber, Krapf, and Schmidheiny
(2022) find reported wealth about 43 percent higher per percentage-point rate
cut in Switzerland, well above our assumption; and Londo\~no-V\'elez and
\'Avila-Mahecha (2021) document that the wealthiest Colombians concealed about
a third of their wealth offshore. These semi-elasticities also serve as a
ceiling that absorbs migration and expatriation, which we do not model
separately, and we recommend publishing sensitivity bands around them.

\textbf{Concealment share of the avoidance response.} Structural assumptions,
set in July 2026, splitting the avoidance response above into concealment (the
asset and its income and estate value leave the reported bases entirely) and
legal valuation gaming (the assessed value is lowballed but the income remains
visible). Marketable assets have observable exchange prices, so their
avoidance is treated as 100 percent concealment; closely held avoidance is
split 50/50, since valuation discounts on private businesses are real and
legal. The concealment reading is consistent with the Colombian evidence,
where top-end evasion took the form of entire assets hidden offshore rather
than shaded valuations.

\textbf{Reported-estate net-of-tax elasticity.} From Kopczuk and Slemrod
(2001), who regress reported estates on the estate tax rates prevailing over
the decedent's life and find an elasticity with respect to the net-of-tax rate
of about 0.16, with pooled estimates ranging from 0.10 to 0.22 across
specifications; we carry the band into our sensitivity analysis. Adopted July
2026 in an exact net-of-tax power form---the reported estate retains the
fraction $((1 - \tau_{\text{scenario}})/(1 - \tau_{\text{baseline}}))^{0.16}$,
evaluated at each record's marginal estate rate---which reproduces the local
elasticity at the current top rate, remains exact for large reforms, and
handles newly taxable estates.

\medskip\noindent\textbf{\small MECHANICAL INTERACTIONS}

\textbf{Share of new tax paid out of saving.} The profile---the share of a persistent above-baseline tax flow paid out of saving rather than
consumption, by age and within-age wealth rank, running from about 0.04 at the
bottom to 0.80 at the top---was calibrated in July 2026. The anchor is the
bridge $s = 1 - \varepsilon \cdot (C/Y)$, with $\varepsilon \approx 0.7$ the
elasticity of consumption to permanent income from Straub (2019) and
consumption-to-income ratios by rank from Mian, Straub, and Sufi (2020): the
top 1 percent's income share near 20 percent against a consumption share of 6
to 7 percent gives $s \approx 0.8$ at the top, while hand-to-mouth behavior
(Kaplan and Violante 2022) pulls the bottom toward zero. The gradient is
cross-checked against the saving-rate gradient in lifetime income in Dynan,
Skinner, and Zeldes (2004), whose top 1 percent saves about half of income
across specifications, and against the liquidity gradient of transitory MPCs
in Fagereng, Holm, and Natvik (2021); the age tilt is attenuated to zero at
top ranks because high-income elderly households do not run down wealth (De
Nardi, French, and Jones 2010). Validation runs in July 2026 put the
dollar-weighted effective share for top-concentrated reforms at about 0.78, at
the top of the predicted band.

\textbf{Corporate incidence constants.} Provisional placeholders pending
direct measurement, and we present corporate-rate results as such. The one
anchored value is the normal-return share of the corporate wedge, 0.375:
Treasury's distribution methodology attributes about 63 percent of the
corporate tax to supernormal returns (Cronin, Lin, Power, and Cooper 2013) and
the Tax Policy Center's about 60 percent (Nunns 2012), implying a normal share
of 0.37 to 0.40; we sweep corners of 0 and 0.5. The remaining constants---the
C-corporation share of the normal capital stock (0.40), the U.S.-taxable
exposure scale, and the equity risk premium among them---carry stated priors
rather than measurements. The normal-return share, the corporate share of
normal capital, and the permanence of the priced-in change are exposed as
sweep switches; the remaining constants are fixed in code.

\subsection*{References}

{\small
\bibentry Agersnap, Ole, and Owen Zidar. 2021. ``The Tax Elasticity of
Capital Gains and Revenue-Maximizing Rates.'' \emph{American Economic Review:
Insights} 3 (4): 399--416.

\bibentry Bakija, Jon, and Bradley T. Heim. 2011. ``How Does Charitable
Giving Respond to Incentives and Income? New Estimates from Panel Data.''
\emph{National Tax Journal} 64 (2, Part 2): 615--650.

\bibentry Br\"ulhart, Marius, Jonathan Gruber, Matthias Krapf, and Kurt
Schmidheiny. 2022. ``Behavioral Responses to Wealth Taxes: Evidence from
Switzerland.'' \emph{American Economic Journal: Economic Policy} 14 (4):
111--150.

\bibentry Cronin, Julie Anne, Emily Y. Lin, Laura Power, and Michael Cooper.
2013. ``Distributing the Corporate Income Tax: Revised U.S. Treasury
Methodology.'' \emph{National Tax Journal} 66 (1): 239--262.

\bibentry DeBacker, Jason, Bradley Heim, and Alexander Yuskavage. 2025.
``Marginal Tax Rates and Income Tax Noncompliance.'' Presented at the
National Tax Association annual meetings, November 2025.

\bibentry De Nardi, Mariacristina, Eric French, and John B. Jones. 2010.
``Why Do the Elderly Save? The Role of Medical Expenses.'' \emph{Journal of
Political Economy} 118 (1): 39--75.

\bibentry Dowd, Tim, Robert McClelland, and Athiphat Muthitacharoen. 2015.
``New Evidence on the Tax Elasticity of Capital Gains.'' \emph{National Tax
Journal} 68 (3): 511--544.

\bibentry Dynan, Karen E., Jonathan Skinner, and Stephen P. Zeldes. 2004.
``Do the Rich Save More?'' \emph{Journal of Political Economy} 112 (2):
397--444.

\bibentry Fagereng, Andreas, Martin B. Holm, and Gisle J. Natvik. 2021.
``MPC Heterogeneity and Household Balance Sheets.'' \emph{American Economic
Journal: Macroeconomics} 13 (4).

\bibentry Jakobsen, Katrine, Kristian Jakobsen, Henrik Kleven, and Gabriel
Zucman. 2020. ``Wealth Taxation and Wealth Accumulation: Theory and Evidence
from Denmark.'' \emph{Quarterly Journal of Economics} 135 (1): 329--388.

\bibentry Kaplan, Greg, and Giovanni L. Violante. 2022. ``The Marginal
Propensity to Consume in Heterogeneous Agent Models.'' \emph{Annual Review of
Economics} 14.

\bibentry Kopczuk, Wojciech, and Joel Slemrod. 2001. ``The Impact of the
Estate Tax on Wealth Accumulation and Avoidance Behavior.'' In
\emph{Rethinking Estate and Gift Taxation}, edited by William G. Gale, James
R. Hines Jr., and Joel Slemrod. Washington, DC: Brookings Institution Press.

\bibentry Londo\~no-V\'elez, Juliana, and Javier \'Avila-Mahecha. 2021.
``Enforcing Wealth Taxes in the Developing World: Quasi-Experimental Evidence
from Colombia.'' \emph{American Economic Review: Insights} 3 (2): 131--148.

\bibentry Mian, Atif, Ludwig Straub, and Amir Sufi. 2020. ``The Saving Glut
of the Rich.'' NBER Working Paper 26941.

\bibentry Nunns, Jim. 2012. ``How TPC Distributes the Corporate Income Tax.''
Urban-Brookings Tax Policy Center.

\bibentry Prisinzano, Richard, and James Pearce. 2018. ``Tax Based Switching
of Business Income.'' Penn Wharton Budget Model Working Paper W2018-2.

\bibentry Randolph, William C. 1995. ``Dynamic Income, Progressive Taxes, and
the Timing of Charitable Contributions.'' \emph{Journal of Political Economy}
103 (4): 709--738.

\bibentry Saez, Emmanuel, Joel Slemrod, and Seth H. Giertz. 2012. ``The
Elasticity of Taxable Income with Respect to Marginal Tax Rates: A Critical
Review.'' \emph{Journal of Economic Literature} 50 (1): 3--50.

\bibentry Seim, David. 2017. ``Behavioral Responses to Wealth Taxes: Evidence
from Sweden.'' \emph{American Economic Journal: Economic Policy} 9 (4):
395--421.

\bibentry Smith, Matthew, Danny Yagan, Owen Zidar, and Eric Zwick. 2019.
``Capitalists in the Twenty-First Century.'' \emph{Quarterly Journal of
Economics} 134 (4): 1675--1745.

\bibentry Straub, Ludwig. 2019. ``Consumption, Savings, and the Distribution
of Permanent Income.'' Working paper, Harvard University.
}

\end{document}
